\documentclass[11pt]{article}
\usepackage[a4paper,margin=27mm]{geometry}
\usepackage{amsmath,amssymb,amsthm,booktabs,hyperref}
\newtheorem{theorem}{Theorem}\newtheorem{lemma}[theorem]{Lemma}
\newcommand{\F}{\mathbb F}\newcommand{\PSL}{\operatorname{PSL}}
\newcommand{\PGL}{\operatorname{PGL}}\newcommand{\SL}{\operatorname{SL}}
\newcommand{\SO}{\operatorname{SO}}
\title{Four homologies in a projective plane}
\author{A recovery of the odd-characteristic rank-four classification in L027}
\date{}
\begin{document}\maketitle
Let $F$ be a finite field of odd characteristic $p$. We classify all rank-four string
C-markings in $\PSL_3(F)$, including the degenerate case, and obtain
the directly regular alternatives needed in Section~6 of
Leemans--Vandenschrick. The normal forms originate in their Section~5.3.
The proof below uses the actual translation kernels and an acquired
projective-line classification, without first classifying all rank-three
subgroups of $\PSL_3(F)$.

We retain the usual string end-intersection criterion, quoted as
Lemma~2.1 of the source, and the classical identification
$\SO_3(F)\cong\PGL_2(F)$. We also acquire the previously proved
Topic E classification: a rank-four string C-marking of $\PSL_2(K)$
or $\PGL_2(K)$ is exactly the $\PSL_2(11)$ marking of type $[3,5,3]$,
the $\PSL_2(19)$ marking of type $[5,3,5]$, or the $\PGL_2(5)$
marking of type $[3,3,3]$. Its finite presentations and complete
intersection checks remain in that acquired component. Dickson's
subgroup theorem, in its actual natural-subfield form, is the only
classification input used in passing to those groups.

\section*{The diagonal outer pair}
\begin{lemma}\label{odd-bound}
A string C-marking in $\PSL_3(F)$ has rank at most four. Every
rank-four marking has, after a projective change of coordinates,
unique order-two lifts of the form
\begin{align*}
 r_0&=\operatorname{diag}(1,-1,-1),&
 r_3&=\operatorname{diag}(-1,-1,1),\\
 r_1&=\begin{pmatrix}a&b&0\\c&-a&0\\0&0&-1\end{pmatrix},&
 r_2&=\begin{pmatrix}-1&0&0\\0&-d&f\\0&e&d\end{pmatrix},
 \qquad bc=1-a^2,\quad ef=1-d^2.
\end{align*}
Here $a,d\notin\{0,-1\}$. If either is one, exactly one of its two
off-diagonal parameters is zero.
\end{lemma}
\begin{proof}
The scalar centre has odd order dividing three, so projective involutions
lift uniquely to involutions of determinant one. Projectively commuting
lifts commute: a scalar commutator $\zeta$ would have $\zeta^2=1$.
In rank at least five, the first, third and fifth generators give an
independent $C_2^3$. Their lifts are simultaneously diagonalizable,
whereas the determinant-one diagonal sign group has order four. This
contradicts independence.

At rank four, simultaneously diagonalize the independent outer pair as
$r_0,r_3$ above. The involution $r_1$ preserves both eigenspaces of
$r_3$. It acts as $-1$ on $\langle e_3\rangle$, since otherwise its
one-dimensional $+1$ eigenspace would be that point and $r_1=r_3$.
Its remaining block has trace zero and determinant $-1$. The same
argument applies to $r_2$, giving the displayed forms.
If $a=0$, then $(r_0r_1)^2=r_3$, contrary to IP. If $a=-1$, then
$(r_0r_1)^p=r_3$, again impossible. If $a=1$ and $b=c=0$, then
$r_1=r_0$. Reversal gives the same assertions for $d,e,f$.
\end{proof}

\section*{The conic alternative}
\begin{lemma}\label{dihedral}
A dihedral group has no nondegenerate rank-three string C-marking.
\end{lemma}
\begin{proof}
All three generators must be reflections: a central involution would
commute with its neighbour. Write the commuting distinct ends as $r$
and $rz$, where $z$ is the central involution. If $s$ is the middle
reflection, then $(s r z)^2=(rs)^{-2}$. Since $|rs|>2$, this is a
nonidentity rotation in both dihedral ends and not in $\langle s\rangle$,
contradicting IP.
\end{proof}

\begin{lemma}\label{conic}
If all four off-diagonal parameters $b,c,e,f$ are nonzero, the marking
is one of the three acquired projective-line markings. The ambient
characteristics are respectively $11,19,5$.
\end{lemma}
\begin{proof}
Direct multiplication shows that the four generators preserve the
nonsingular diagonal quadratic form
\[
 Q=(c/b)x_1^2+x_2^2+(f/e)x_3^2.
\]
The resulting orthogonal group has trivial intersection with the scalar
centre of $\SL_3(F)$, since a scalar preserving $Q$ has square one.
Thus the marked group embeds faithfully in $\PGL_2(F)$. Its outer
labels exceed two by Lemma~\ref{odd-bound}, and the middle generators
do not commute: the $(1,3)$ entry of $r_1r_2$ is $bf\ne0$, whereas
that of $r_2r_1$ is zero. The marking is nondegenerate.

We apply Dickson's subgroup alternatives to this actual subgroup.
Cyclic groups and the involution-generated subgroup $V_4$ of $A_4$
are excluded. In an odd-characteristic affine group every involution
has the form $x\mapsto-x+v$, and two distinct such involutions do not
commute; the commuting first and third generators exclude that case.
The dihedral case is excluded by Lemma~\ref{dihedral} applied to an
end. If the whole group were $S_4$ or $A_5$, either proper rank-three
end would lie in an $A_4$, a dihedral subgroup or an abelian subgroup,
by their elementary proper-subgroup lists; none admits a nondegenerate
rank-three marking. IP makes these ends proper, so these two cases
are excluded as well.

The group is therefore a natural subfield $\PSL_2(K)$ or $\PGL_2(K)$.
The acquired Topic E theorem gives the three marked possibilities.
The orders and simplicity distinguish them: $\PSL_2(11)$ and
$\PSL_2(19)$ are simple, while $\PGL_2(5)$ has an index-two subgroup.
Their natural subfield embeddings consequently have characteristic
$11,19,5$, respectively. Conversely their standard three-dimensional
orthogonal representations embed them in $\PSL_3(F)$ in those
characteristics, with the same actual markings.
\end{proof}

\section*{The affine alternatives}
Duality transposes each lift, interchanging $b,c$ and also $e,f$.
Reversal, followed by interchange of the first and third coordinates,
interchanges the two pairs. These operations reduce all remaining cases
to $b=0$, with either no further zero, $e=0$, or $f=0$.
The nonzero parameter $c$ may be scaled to one; when $d=1$, the other
nonzero off-diagonal parameter can also be scaled to one. These are
changes of basis, not changes to the ordered generated group.

\begin{lemma}\label{one-zero}
The case $b=0$ and $ef\ne0$ cannot have IP.
\end{lemma}
\begin{proof}
Here $a=1$, $c\ne0$, and $d\notin\{0,1,-1\}$. In block form with
the invariant plane $\langle e_2,e_3\rangle$, write
\[
 r_1=\begin{pmatrix}1&0\\c e_1&-I_2\end{pmatrix},\quad
 r_2=\operatorname{diag}(-1,P),\quad
 r_3=\operatorname{diag}(-1,Q_0).
\]
Here
\[
 P=\begin{pmatrix}-d&f\\e&d\end{pmatrix},\qquad
 Q_0=\operatorname{diag}(-1,1).
\]
Let $T(v)=\left(\begin{smallmatrix}1&0\\v&I_2\end{smallmatrix}\right)$.
The actual subgroup $B=\langle r_1,r_2,r_3\rangle$ contains
$(r_1r_2)^2=T((1-d)c,ec)$. Its intersection with the additive
translation group is invariant under $-P$ and $-Q_0$.
The latter is the sign change of the second coordinate. Adding and
subtracting its image and dividing by two projects onto each coordinate.
Hence this intersection contains $T((1-d)c,0)$. Applying $-P$ and
projecting again multiplies that coordinate by $d$. Since
$\F_p[d]$ is a field and $1-d\ne0$, its polynomial span contains
$T(c,0)$. But $r_1=T(c,0)r_0$, so $r_0\in B$, contrary to IP.
\end{proof}

\begin{theorem}\label{affine}
The two remaining affine cases are as follows.
\begin{enumerate}
\item If $b=e=0$, the marking has type $[p,2,p]$ and group
$D_{2p}\times D_{2p}$.
\item If $b=f=0$, the marking has type $[p,2p,p]$ and group
$\operatorname{Heis}(p)\rtimes C_2^2$, where the Heisenberg group
consists of all lower unitriangular $3\times3$ matrices over $\F_p$.
\end{enumerate}
Each displayed normal form has full IP and is directly regular.
\end{theorem}
\begin{proof}
In the first case, after scaling $c=f=1$, the two cyclic translation
groups $\langle I+E_{21}\rangle$ and $\langle I+E_{23}\rangle$
commute. The diagonal $r_0$ inverts the first and fixes the second;
$r_3$ does the reverse. This is the actual direct product of two
dihedral groups, with its two rank-two markings placed consecutively.
The type and all subset intersections follow from that direct product.

In the second case set $c=e=1$. Put
\[
 h(x,y,z)=\begin{pmatrix}1&0&0\\x&1&0\\z&y&1\end{pmatrix},
 \quad D_0=r_0,\quad D_3=r_3.
\]
The root groups $h(x,0,0)$ and $h(0,y,0)$, generated respectively
by $r_0r_1$ and $r_2r_3$, have commutators giving all $h(0,0,z)$,
with $x,y,z\in\F_p$. The two diagonal generators normalize this
Heisenberg group and intersect it trivially. The full group thus consists
uniquely of $h(x,y,z)D_0^iD_3^j$, with $i,j\in\{0,1\}$, and has
order $4p^3$. Here
\[
 r_1=h(1,0,0)D_0,\qquad r_2=h(0,-1,0)D_3.
\]
The first actual end $A=\langle r_0,r_1,r_2\rangle$ consists of
these normal forms with $y=-j$, and the second
$B=\langle r_1,r_2,r_3\rangle$ consists of those with $x=i$.
For example $A$ contains the full $x,z$ translation plane (conjugate
the $x$ root group by $r_2$), and its complement is the commuting
pair $r_0,r_2$. This proves both inclusions in its displayed description;
the other end follows in the same way.

Their intersection has $x=i,y=-j$ and arbitrary $z$, hence order $4p$.
Direct multiplication gives
$(r_1r_2)^2=h(0,0,1)$ and shows that $r_1r_2$ has order $2p$.
Consequently $A\cap B=\langle r_1,r_2\rangle$. Within $A$,
the subgroup $\langle r_0,r_1\rangle$ has $j=y=z=0$;
intersecting it with the preceding description of
$\langle r_1,r_2\rangle$ leaves exactly $\langle r_1\rangle$.
Reversal gives the other rank-three test. The end criterion proves
full IP. The outer products have order $p$, giving type $[p,2p,p]$.

Finally, in both cases the map to $C_2$ which sends a diagonal
$D_0^iD_3^j$ to $i+j$ sends all four generators to its nonidentity
element. Its kernel is exactly the even-word subgroup, proving direct
regularity.
\end{proof}

\enlargethispage{2\baselineskip}
\section*{Classification and source qualifications}
Combining the lemmas gives exactly the source's Theorem~5.3: the
degenerate direct product, the Heisenberg extension, and the three conic
markings. The nondegenerate directly regular possibilities are only the
Heisenberg marking of type $[p,2p,p]$ and the simplex marking of type
$[3,3,3]$ in characteristic five. The two simple conic groups have no
index-two subgroup and are not directly regular. The simplex marking
is directly regular by the usual sign map on $S_5$.

The source's physical page~17 names $D_{2p}$ in the two-zero direct
product case; its Theorem~5.3 correctly gives $D_{2p}^2$. The actual
matrix decomposition above settles this discrepancy. Its two-point
description also requires allowing duality: our chosen form fixes two
lines, and its dual fixes two points. The restriction
``$q=5$'' in the rank-six discussion must mean characteristic five
for this subgroup alternative, since the same natural subfield marking
embeds for every extension of $\F_5$. This is not yet a proof of the
rank-six exclusion. The source's Table~1 parameter calculations are
replaced here by the acquired complete marking classification; no new
finite census is charged to that previously recovered result. The broader
rank-three classifications in Theorems~5.1--5.2 remain a separate
source-contract item.
\end{document}
